\chapter{Solutions to Section B}\label{ch:soltracing}

Each box is the real output of the program, captured when the book was built.

\solution{Tracing}{1}
\outputfile{t01}
The three appends succeed: $r=3$. Both front insertions succeed, $r=1\cdot10+1=11$. The
list is now full, so \cc{append(\&l,42)} returns 0. \cc{removeLast} only decrements
\cc{size}; \cc{removeFirst} shifts left, leaving the raw array \cc{[16 4 8 8 15]} with
size 3. \cc{setItem(\&l,99,3)} is valid; position 4 exceeds the size and fails.
\emph{Twist:} \cc{itemAt} does not check its position, so \cc{itemAt(l,4)} reads the
stale cell \cc{items[3]}, which still holds 8. After \cc{append(\&l,5)} the same cell
holds 5.

\solution{Tracing}{2}
\outputfile{t02}
A: the struct (with its array) was copied; the caller sees nothing. B: through the
pointer, both fields change. C: \cc{a.items} decays to a pointer to the caller's array,
so \cc{items[0]} changes but \cc{size} cannot. D: the same for a plain array.
E: \cc{llByValue} receives a copy of \cc{head} and \cc{size}; resetting them affects only
the copy, but \cc{l.head->item = 100} writes into the \emph{shared} first node.

\solution{Tracing}{3}
\outputfile{t03}
$c_1=n^2=256$. $c_2=1+2+\dots+n=n(n+1)/2=136$. $c_3$: the outer loop runs for
$i=1,2,4,8$, that is $\lg n=4$ times, so $4\cdot16=64=n\lg n$. $c_4$: $i=16,8,4,2,1$, five
times $=\lg n+1$. $c_5$: for each $i$ the inner loop runs $\lceil(16-i)/4\rceil$ times;
$4(4+3+2+1)=40$, which is $\Theta(n^2)$ (about $n^2/8$). $c_6=1+2+4+8=15=n-1$:
\emph{twist} --- a doubling outer loop whose inner loop runs $i$ times is $\Theta(n)$,
not $\Theta(n\log n)$.

\solution{Tracing}{4}
\outputfile{t04}
\cc{erase(1)} copies each item one cell left but does not clear the old last cell, so
\cc{60} appears twice in the raw array. \cc{insert(3,99)} overwrites that stale copy
when shifting right. \cc{erase(5)} removes the real 60 ($r=1$); the second
\cc{erase(5)} fails because now $\cc{size}=5$ ($r=1+0=1$). The raw array is unchanged by
the size decrement. Two front insertions shift everything. \emph{Twist:} both later
inserts fail --- \cc{insert(9,1)} because $9>\cc{size}$, and \cc{insert(8,1)} because
$8>7=\cc{size}$, even though the array still has one free cell.

\solution{Tracing}{5}
\outputfile{t05}
The removal only advances \cc{j}; the 1 stays in cell 0. Adding 4 and 5 fills cells 3
and then 0 (wrap-around). The queue is then full and \cc{add(\&q,6)} returns 0. Two
removals take 2 and 3; \cc{j} becomes 3. Then 7 and 8 go to cells $(3+2)\bmod4=1$ and
$2$. The final loop prints the queue in queue order (4 5 7 8), not array order.
\emph{Twist:} after the queue is emptied, \cc{show} still prints the old values; only
\cc{n} says the queue is empty.

\solution{Tracing}{6}
\outputfile{t06}
\cc{popBack} walks from the head to the node before the tail; the second \cc{popBack}
leaves one node, which is both head and tail. \cc{popFront} on the one-node list must
reset \emph{both} pointers to \cc{NULL}. \emph{Twist:} the failed \cc{popFront} returns 0
and leaves \cc{x} at the last value successfully popped, 3.

\solution{Tracing}{7}
\outputfile{t07}
\cc{appendA} receives a copy of the head pointer. On an empty list it stores the new node
in that copy only: \cc{a} stays \cc{NULL} and the node leaks, three times. \cc{appendB}
receives the address of the caller's pointer and works. \emph{Twist:} once the list is
non-empty (\cc{c} starts with 7), \cc{appendA} works as well, since it modifies the
\cc{next} field of a shared node: both lists have length 4.

\solution{Tracing}{8}
\outputfile{t08}
Removing all 2's removes the first two nodes (the head changes twice), the one in the
middle, and the last two (the tail changes twice); \cc{5 7} remain and the backward
display shows that the \cc{prev} links are right. Erasing a missing value changes nothing.
After removing 7, the single node 5 is both head and tail and its \cc{prev} is
\cc{NULL}. Removing it empties the list: both pointers are \cc{NULL}.

\solution{Tracing}{9}
\outputfile{t09}
\cc{mystery} prints each item on the way down, \cc{\#} at the end, and each item again
on the way back: the list forwards and then backwards. \cc{g} adds the even items and
subtracts 1 for every odd item: $2+4-2=4$ and $2+8+6-3=13$. \cc{h} counts the positions
where an item is smaller than its successor: $1<2<3<4$ gives 3; in
$3,5,2,8,6,9$ the ascents are $3<5$, $2<8$, $6<9$: also 3.

\solution{Tracing}{10}
\outputfile{t10}
Line 1: the stack gives its newest item 3, the queue its oldest 10. Then 3 joins the
queue (\cc{[20 3]}) and 10 the stack (\cc{[1 2 10]}). Line 2: the loop pops until the stack is
empty, three times (10, 2, 1): $r=3$ and \cc{x} is the last popped item, 1. Line 3: two
dequeues (20, 3): $r=2$, \cc{y}$=3$. Last line: 4 and 5 go stack $\to$ queue (\cc{5 4}) $\to$ stack
(top 4). \emph{Twist:} two reversals cancel, so the items come out in push order.

\solution{Tracing}{11}
\outputfile{t11}
The stack grows to \cc{+ * ( -} (depth 4) inside the parentheses. At \cc{)} the \cc{-} is
output and \cc{(} discarded. At \cc{/} the \cc{*} has equal precedence and leaves first.
\emph{Twist:} \cc{prec('(')} is 0, so the operator loop stops at a parenthesis without an
explicit test.

\solution{Tracing}{12}
\outputfile{t12}
The first dequeue moves 1, 2, 3 to \cc{out}. The dequeues of 2 and 3 need no transfer,
although 4, 5, 6 are waiting in \cc{in}. Only when \cc{out} is empty again are 4, 5, 6
moved. In total each item is moved once: \emph{6} moves for 6 items. The last dequeue
finds both stacks empty and returns $-1$.

\solution{Tracing}{13}
\outputfile{t13}
\cc{evalB} computes $b\ \mathrm{op}\ a$ with the operands swapped. For \cc{+} and
\cc{*} the error is invisible. For \cc{93-} it gives $3-9=-6$; for \cc{82/} it gives
$2/8=0$. In \cc{27-3/}, \cc{evalA} computes $(2-7)/3=-5/3=-1$ (truncation toward
zero, not $-2$), while \cc{evalB} computes $3/(7-2)=0$. The errors compound: for
\cc{92-3-}, \cc{evalB} gives $3-(2-9)=10$.

\solution{Tracing}{14}
\outputfile{t14}
\cc{3 2 1 4}: push three, pop three, push, pop. \cc{1 2 3 4}: alternate. \cc{3 1 2 4}:
after 3 is printed, 2 is on top and 1 is below it, so the function pushes 4 (one more
\cc{u}) and then has nothing left to push: stuck with top 4. This is the forbidden
pattern $3,1,2$. \cc{2 4 3 1} works. \cc{4 1 2 3}: after 4 the top is 3 and 1 is
buried. Three of the five are stack permutations.

\solution{Tracing}{15}
\outputfile{t15}
70 absorbs the 60 before it (span $1+1$). 75 absorbs \cc{60:1} and \cc{70:2}: span
$1+1+2=4$. 85 absorbs \cc{75:4} and \cc{80:1}: span 6. The second 85 absorbs the first
(\cc{<=}): span 7. \emph{Twist:} the stack keeps only 2 pairs at the end although 8 days
were seen; the absorbed days are represented by the spans.

\solution{Tracing}{16}
\outputfile{t16}
The middle is index $\lfloor(n-1)/2\rfloor$: for $n=1..4$ that is $0,0,1,1$. A push at
the back moves \cc{mid} forward when $n$ becomes odd; a push at the front moves it
\emph{backward} when $n$ becomes even (from 20 to 10 at $n=6$). Pops do the reverse.

\solution{Tracing}{17}
\outputfile{t17}
For $x=4$ the search keeps going left after finding a 4, ending at the \emph{first} 4
(index 1). For $x=10$ the answer 6 is where 10 would be inserted. For $x=20$ all items are
smaller and the answer is $n=9$, which is not a valid index: a caller must test
$\cc{answer}<n$ before reading \cc{a[answer]}. Each search takes at most
$\lceil\lg(n+1)\rceil=4$ steps.

\solution{Tracing}{18}
\outputfile{t18}
Queue after each round: \cc{[4 5 6 1 2]} (3 out), \cc{[1 2 4 5]} (6 out),
\cc{[5 1 2]} (4 out), \cc{[5 1]} (2 out); in the last round 5 and 1 both go to the back
and the third person counted is 5 again, so 5 leaves. The winner is 1. Recurrence with $k=3$: $J=0,1,1,0,3,0$ for
$i=1..6$, so the winner is $J(6)+1=1$. \emph{Twist:} when fewer than $k$ people remain,
the count wraps around the circle more than once.

\solution{Tracing}{19}
\outputfile{t19}
\cc{splice(\&a,2,\&b)} makes 10 the item at index 2. \cc{a.head->next->next} is the node
10, so \cc{c} goes after it. Splicing at the end must update \cc{a.tail} (the tail
stays 5 in the middle splices). An invalid position changes nothing, and \cc{d} keeps
its node until the valid splice at \cc{a.size} appends it and moves the tail.

\solution{Tracing}{20}
\outputfile{t20}
\cc{sh} shares its nodes with \cc{orig}: changing \cc{sh.head->item} changes
\cc{orig}'s first item to 99. \cc{dp} has its own nodes. \emph{Twist:}
\cc{pushFront(\&sh,0)} changes only \cc{sh.head} and \cc{sh.size}: \cc{orig} does not see
the new node and still has size 3, although the two lists share their last three nodes.

\solution{Tracing}{21}
\outputfile{t21}
\cc{pushBack(1)} writes cell 0. \cc{pushFront(2)}: $f=(0+4)\bmod5=4$. \cc{pushFront(3)}:
$f=3$. \cc{pushBack(4)} writes cell $(3+3)\bmod5=1$. \cc{popBack} returns cell
$(3+3)\bmod5=1$, i.e.\ 4; \cc{popFront} returns 3. Then 5 and 6 fill cells 1 and 2 and
7 goes to cell 3; the deque is full and both further pushes fail ($r=0$). The last line
shows why the code adds 4: \cc{(0-1)\%5} is $-1$ in C, \cc{(0-1+5)\%5} is 4, and
\cc{-7/2} is $-3$.

\solution{Tracing}{22}
\outputfile{t22}
\cc{add(1,15)} with $n=4$: $1<2$, so the left side (item 10) moves one cell left, to cell
7, and $j=7$. \cc{add(3,25)} with $n=5$: $3\ge2$, so items 3 and 4 (30, 40) move right.
\cc{add(0,5)}: nothing moves, $j$ steps back to 6. \cc{add(6,45)} with $n=7$: only the
last item moves. Total moves 4, against $3+2+6+1=12$ for an array list doing the same
insertions.

\solution{Tracing}{23}
\outputfile{t23}
At $i=2$ the 12 pops both earlier indices. At $i=4$ the 8 pops the 3. At $i=5$ the front
index 2 has left the window ($2\le5-3$) and is popped from the front; then the new 8 pops
the old 8 because the test is \cc{<=}, so the deque keeps the \emph{newer} of two equal
values, which stays in the window longer. The maxima are 12, 12, 12, 8, 8, 8.

\solution{Tracing}{24}
\outputfile{t24}
The list is a circle, so moving \cc{last} one step \emph{rotates} the queue: the old front
1 becomes the rear. Two more steps bring it back. \cc{deq} removes the front 1. After the
loop empties the queue, \cc{last} is \cc{NULL}. \emph{Twist:} the final failed \cc{deq}
returns 0 and \cc{x} keeps 4, the last item removed.

\solution{Tracing}{25}
\outputfile{t25}
The empty deque's sentinel points to itself and the size is 0.
After \cc{0 1 2 3}, reading and removing both ends gives 3 and 0. The last item is 2 and
the first node's \cc{prev} is the sentinel. \emph{Twist:} \cc{LinkedDeque e = d;} copies
only the members, the \emph{pointer} \cc{sentinel} and \cc{count}: both objects use the
same sentinel (1) and \cc{e.size()} is 2. After \cc{d.push\_front(9)} the walk from
\cc{e} also starts at 9 (shared nodes), but \cc{e.size()} is still 2. After
\cc{d.clear()} the shared list is empty (the last 1), yet \cc{e.size()} still says 2: a
shallow copy goes out of date.
